\documentclass[10pt,aspectratio=169]{beamer}

\usetheme{Madrid}
\usecolortheme{default}
\setbeamertemplate{navigation symbols}{}

\usepackage{amsmath}
\usepackage{booktabs}
\usepackage{graphicx}
\usepackage{tikz}
\usetikzlibrary{positioning,arrows.meta}

% Notes setup: keep hidden by default.
% For presenter notes output, uncomment one of:
% \setbeameroption{show notes}
% \setbeameroption{show notes on second screen=right}
\setbeameroption{hide notes}

\title[RL for Octopus Memory Pooling]{RL-Based Memory Pooling for Octopus CXL Systems}
\subtitle{8-minute advisor update}
\author{Pulak Mehrotra}
\date{\today}

% Figure helper for fast iteration when files move.
\newcommand{\figorplaceholder}[2]{%
  \IfFileExists{#1}{%
    \includegraphics[width=\linewidth,height=#2,keepaspectratio]{#1}%
  }{%
    \fbox{\parbox[c][#2][c]{0.9\linewidth}{\centering Missing figure:\\\texttt{#1}}}%
  }%
}

\begin{document}

% ----------------------------------------------------------------
\begin{frame}
  \titlepage
  \note{
  Talk track: framing for an 8-minute technical overview.
  Main message to advisor: RL is not yet better than greedy everywhere, but with
  temporal signals and better shaping it is moving toward the gap to optimal.
  }
\end{frame}

% ----------------------------------------------------------------
\begin{frame}{1) What is RL?}
  \begin{itemize}
    \item Reinforcement Learning (RL) learns a policy by trial-and-error interaction.
    \item At each decision event:
      \[
      s_t \xrightarrow{\pi_\theta} a_t,\quad r_t,\quad s_{t+1}
      \]
    \item Objective:
      \[
      \max_{\pi_\theta}\ \mathbb{E}\!\left[\sum_{t=0}^{T}\gamma^t r_t\right]
      \]
    \item In this project, each VM arrival is one decision point.
    \item The policy learns how to split VM memory demand across reachable MPDs.
  \end{itemize}
  \note{
  Keep this concise and conceptual. Do not dive into SAC internals yet.
  Emphasize event-driven decision making (VM arrivals), not fixed clock steps.
  }
\end{frame}

% ----------------------------------------------------------------
\begin{frame}{2) Why RL is a good fit (temporal trends)}
  \figorplaceholder{../figs/demand_overview.pdf}{0.72\textheight}
  \vspace{0.3em}
  \begin{itemize}
    \item Clear diurnal regularity means future demand is partially predictable.
    \item Time-aware allocation can reserve headroom before expected ramps.
    \item RL is well-suited to learning this anticipation from trace replay.
  \end{itemize}
  \note{
  Make the figure the center of this slide.
  Keep the narrative focused on temporal predictability and anticipation.
  }
\end{frame}

% ----------------------------------------------------------------
\begin{frame}{3) Policy Design: State}
  \textbf{State table (cleaned-up design)}
  \scriptsize
  \begin{tabular}{p{0.30\linewidth}p{0.34\linewidth}p{0.30\linewidth}}
    \toprule
    \textbf{Signal} & \textbf{Current implementation} & \textbf{Next iteration} \\
    \midrule
    Per-accessible-MPD load & Normalized current loads & Add short-horizon deltas/trends \\
    Topology connectivity & Binary reachable-MPD mask & Keep (hard validity constraint) \\
    Request context & Normalized VM memory request & Keep and add request-type bucket \\
    Global pressure & Current peak normalized MPD load & Keep and add rolling peak estimate \\
    Temporal context & Hour-of-day $\sin/\cos$ & Keep and add day-of-week encoding \\
    \bottomrule
  \end{tabular}
  \normalsize
  \vspace{0.3em}
  \textbf{Difference vs current implementation in \texttt{train.py}/env}
  \begin{itemize}
    \item Current code is mostly instantaneous.
    \item Next iteration should add short history/trend features (e.g., load deltas)
      to improve anticipation.
  \end{itemize}
  \note{
  Be explicit that this slide is the clean conceptual design.
  Mention current env already includes hour sin/cos + load/mask/request/peak.
  Clarify planned extension: trend features for better temporal forecasting.
  }
\end{frame}

% ----------------------------------------------------------------
\begin{frame}{4) Policy Design: Action}
  \begin{itemize}
    \item Output is a continuous vector over reachable MPDs:
      \[
      \mathbf{w} \in \mathbb{R}^{k},\quad w_i \ge 0,\quad \sum_i w_i=1
      \]
    \item Connectivity mask enforces invalid links $\rightarrow$ zero allocation.
    \item Physical allocation:
      \[
      \text{alloc}_i = w_i \cdot \text{VM\_memory}
      \]
    \item This keeps control high-level and low-overhead: no migration action in-loop.
  \end{itemize}
  \vspace{0.4em}
  \textbf{Difference vs current code}
  \begin{itemize}
    \item Current policy outputs unconstrained scores then applies softmax.
    \item We keep this but emphasize explicit masked simplex interpretation.
  \end{itemize}
  \note{
  Advisor may ask: "why not discrete argmin action?" Answer: continuous split is
  expressive and fits SAC naturally.
  }
\end{frame}

% ----------------------------------------------------------------
\begin{frame}{5) Policy Design: Reward}
  \textbf{Primary objective: reduce required per-MPD peak capacity}
  \[
  \text{Savings} = 1 - \frac{m \cdot \max_j(\text{PeakLoad}_j)}{\sum_i \text{DRAM}_i}
  \]
  \vspace{0.4em}
  \textbf{Training reward (clean form)}
  \[
  r_t = -\Delta\max_j(c_j) - \lambda \cdot \mathrm{Var}(\mathbf{c}_t)
  \]
  where $\mathbf{c}_t$ is MPD load vector.
  \begin{itemize}
    \item First term penalizes increasing the worst MPD.
    \item Second term gives dense signal for load balancing.
  \end{itemize}
  \vspace{0.4em}
  \textbf{Difference vs current \texttt{train.py} environment}
  \begin{itemize}
    \item Current reward is close to peak-change only.
    \item Planned update adds variance shaping to reduce mode collapse.
  \end{itemize}
  \note{
  Important bridge: "we are getting there with RL" because reward shaping can
  close observed imbalance gap.
  }
\end{frame}

% ----------------------------------------------------------------
\begin{frame}{6) Model Choice: Why SAC (and where PPO fits)}
  \begin{itemize}
    \item \textbf{Chosen primary model: SAC}
      \begin{itemize}
        \item Continuous action support fits allocation vector directly.
        \item Off-policy replay improves sample efficiency on long traces.
        \item Entropy term helps exploration and avoids brittle deterministic policy.
      \end{itemize}
    \item \textbf{PPO as secondary baseline}
      \begin{itemize}
        \item On-policy, often more stable but less sample efficient here.
      \end{itemize}
    \item \textbf{Practical note}
      \begin{itemize}
        \item Current \texttt{train.py} supports both SAC/PPO; experiments currently centered on SAC.
      \end{itemize}
  \end{itemize}
  \note{
  Keep to one sentence each for SAC/PPO tradeoff.
  If asked: "why not model-based?" answer with engineering complexity + trace-driven setup.
  }
\end{frame}

% ----------------------------------------------------------------
\begin{frame}{7) Results: Where RL stands now}
  \begin{columns}[T,onlytextwidth]
    \column{0.68\textwidth}
      \figorplaceholder{../figs/3panel_by_cxl.pdf}{0.74\textheight}
    \column{0.32\textwidth}
      \textbf{Current status}
      \begin{itemize}
        \item RL is improving but not consistently closing greedy$\rightarrow$optimal gap.
        \item Best gains appear in temporally structured regimes.
      \end{itemize}
      \vspace{0.3em}
      \textbf{Message for this update}
      \begin{itemize}
        \item \textbf{We are getting there with RL.}
        \item Next step: reward/state refinement + longer training for robust gains.
      \end{itemize}
  \end{columns}
  \note{
  Start results with the RL-centered view per advisor presentation flow.
  }
\end{frame}

% ----------------------------------------------------------------
\begin{frame}{8) Results: Load balancing behavior (Greedy vs RL)}
  \begin{columns}[T,onlytextwidth]
    \column{0.5\textwidth}
      \figorplaceholder{../figs/mhd_timeseries_greedy.png}{0.60\textheight}
      \centering \small Greedy
    \column{0.5\textwidth}
      \figorplaceholder{../figs/mhd_timeseries_rl.png}{0.60\textheight}
      \centering \small RL (SAC)
  \end{columns}
  \vspace{0.35em}
  \begin{block}{Takeaway}
    RL can still over-concentrate load on a few MPDs; improving reward shaping
    and temporal features is key to closing the remaining performance gap.
  \end{block}
  \note{
  Use this slide to explain *why* RL is not yet consistently better:
  behaviorally, it sometimes learns imbalance.
  Keep the takeaway short and constructive.
  }
\end{frame}

% ----------------------------------------------------------------
\begin{frame}{9) Results: Topology-level trend}
  \begin{columns}[T,onlytextwidth]
    \column{0.68\textwidth}
      \figorplaceholder{../figs/fc_vs_octopus.pdf}{0.74\textheight}
    \column{0.32\textwidth}
      \textbf{Reading this plot}
      \begin{itemize}
        \item Fully-connected remains an upper benchmark.
        \item Sparse Octopus is harder due to constrained routing.
        \item Gap widens at higher CXL fraction.
      \end{itemize}
      \vspace{0.4em}
      \textbf{Takeaway}
      \begin{itemize}
        \item There is measurable headroom to recover via better policies.
      \end{itemize}
  \end{columns}
  \note{
  Use this as context for why the RL improvements still matter.
  }
\end{frame}

% ----------------------------------------------------------------
\appendix

\begin{frame}{Appendix A: Diurnal trends (backup)}
  \figorplaceholder{../figs/demand_diurnal.pdf}{0.78\textheight}
  \note{
  Backup version of Slide 2 for deeper discussion.
  }
\end{frame}

\begin{frame}{Appendix B: Additional temporal evidence}
  \figorplaceholder{../figs/demand_overview.pdf}{0.78\textheight}
  \note{
  Backup for "why temporal structure exists?" questions.
  }
\end{frame}

\begin{frame}{Appendix C: Peak-to-mean context}
  \figorplaceholder{../figs/peak_to_mean.pdf}{0.78\textheight}
  \note{
  Backup for "is there enough headroom to matter economically?" questions.
  }
\end{frame}

\begin{frame}{Appendix D: Load-balancing behavior details}
  \begin{columns}[T,onlytextwidth]
    \column{0.5\textwidth}
      \figorplaceholder{../figs/mhd_timeseries_greedy.png}{0.68\textheight}
    \column{0.5\textwidth}
      \figorplaceholder{../figs/mhd_timeseries_rl.png}{0.68\textheight}
  \end{columns}
  \vspace{0.2em}
  \small Useful for discussing where RL still diverges from balanced allocation.
  \note{
  This is a strong diagnostic slide. Use only if advisor asks why RL underperforms
  in some settings.
  }
\end{frame}

\end{document}

